% ============================
% §1 数学包
% ============================
\usepackage{amsmath}                                    % 排版数学公式
\usepackage{latexsym}
\usepackage{amsfonts}                                   % AMS 数学字体（Computer Modern 补充符号字体）
\usepackage{amssymb}                                    % AMS 数学符号命令
\usepackage{amscd}                                      % 交换图环境（如不用可删除）
\usepackage{mathrsfs}                                   % 数学 RSFS 书写字体
\usepackage{bm}                                         % 数学黑体（\bm 命令）
\usepackage{physics}                                    % 物理公式排版
\usepackage[math-style=ISO]{unicode-math}               % 使用 ISO 数学风格

% 计数器依赖 section
\numberwithin{equation}{section}
\numberwithin{table}{section}
\numberwithin{figure}{section}

% ============================
% §2 排版格式包
% ============================
\usepackage{graphicx}                                   % 支持插图
\usepackage{color,xcolor}                               % 支持彩色
\usepackage{titlesec}                                   % 设置章节格式
\usepackage{tcolorbox}                                  % 彩色文本框
\usepackage{enumerate}                                  % 更改 enumerate 环境格式
\usepackage[perpage]{footmisc}                          % 脚注每页重新编号

% 设置 paragraph 为只有一级编号（从1开始）
\renewcommand{\theparagraph}{\arabic{paragraph}}
\makeatletter
\@addtoreset{paragraph}{section}                        % 每节开始时重置 paragraph 计数器
\makeatother

% 设置 paragraph 为 block 风格，自动换行
\titleformat{\paragraph}[block]
  {\normalfont\normalsize\bfseries}{\theparagraph.}{0.5em}{}

% 设置标题后换行（after-sep 设为 \newline）
\titlespacing*{\paragraph}{0pt}{3.25ex plus 1ex minus .2ex}{0pt}

% ============================
% §3 代码与算法包
% ============================
\usepackage{minted}                                     % 代码高亮
\usepackage[linesnumbered,ruled,vlined]{algorithm2e}    % 算法排版

% ============================
% §4 表格与图表包
% ============================
\usepackage{diagbox}                                    % 表格斜线表头
\usepackage{tabularx}                                   % 自动调整列宽
\usepackage{multirow}                                   % 合并行
\usepackage{subcaption}                                 % 子图和子标题
\usepackage{minipage-marginpar}                         % minipage 中使用 marginpar
\usepackage{float}                                      % 提供 float 浮动环境
\usepackage{booktabs}                                   % 三线表（\toprule、\midrule、\bottomrule）
\usepackage{listings}                                   % 代码排版（非高亮）

% ============================
% §5 引用与超链接包
% ============================
\usepackage{hyperref}

% 改变超链接颜色
\hypersetup{
    colorlinks=true,
    linkcolor=blue,
    filecolor=blue,
    urlcolor=blue,
    citecolor=cyan,
}

% autoref 中文引用名
\def\equationautorefname{式}
\def\footnoteautorefname{脚注}
\def\itemautorefname{项}
\def\figureautorefname{图}
\def\tableautorefname{表}
\def\partautorefname{篇}
\def\appendixautorefname{附录}
\def\chapterautorefname{章}
\def\sectionautorefname{节}
\def\subsectionautorefname{小小节}
\def\subsubsectionautorefname{小小小节}
\def\paragraphautorefname{段落}
\def\subparagraphautorefname{子段落}
\def\FancyVerbLineautorefname{行}
\def\theoremautorefname{定理}

% ============================
% §6 标题层级与目录设置
% ============================
\setcounter{tocdepth}{3}
\setcounter{secnumdepth}{4}

% ============================
% §7 自定义命令
% ============================
% 角标引用文献和交叉引用
\newcommand{\scite}[1]{\textsuperscript{\cite{#1}}}
\newcommand{\sref}[1]{\textsuperscript{\ref{#1}}}
% boldsymbol 简写别名（等价于 bm 包的 \bm）
\newcommand{\bs}[1]{\boldsymbol{#1}}
% 罗马数字
\newcommand{\romannumber}[1]{\uppercase\expandafter{\romannumeral#1}}
% 红色警示文本
\newcommand{\alert}[1]{\textcolor{red}{#1}}

% 彩色文本框，宽度自适应版心
\newcommand{\tbox}[1]{
  \begin{center}
    \begin{tcolorbox}[colback=gray!10,
        colframe=black,
        width=\linewidth,
        arc=1mm, auto outer arc,
        boxrule=0.5pt,
      ]
      {#1}
    \end{tcolorbox}
  \end{center}
}

% 法语强调下划线
\newcommand{\fri}[1]{\underline{#1}}
